\subsection{Networking, Volumes and Configuration}

\textbf{Single bridge network.} All containers attach to the same Docker bridge network, \texttt{hospital-net}. Traefik resolves service names against Docker's internal DNS, so the rules need nothing more than the service name and the published port. The only entries published to the host machine are port 80 (Traefik HTTP), port 8090 (Traefik dashboard), ports 5672 and 15673 (RabbitMQ AMQP and management), and port 5432 on \texttt{clinical-db} (only for the ETL profile).

\textbf{Host-based routing.} Traefik inspects the \texttt{Host} header: \texttt{auth.localhost}, \texttt{users.localhost}, \texttt{clinical.localhost}, \texttt{analytics.localhost}, \texttt{ml.localhost}, \texttt{mail.localhost} and \texttt{queue.localhost} each map to the corresponding back-end service. The \texttt{.localhost} suffix resolves to \texttt{127.0.0.1} on every modern operating system, so developers reach the services through their browser without modifying \texttt{/etc/hosts}.

\textbf{Configuration contract.} Each service consumes its configuration through environment variables, never through committed files: \texttt{POSTGRES\_USERNAME}, \texttt{POSTGRES\_PASSWORD}, \texttt{POSTGRES\_HOST}, \texttt{POSTGRES\_PORT}, \texttt{POSTGRES\_DB} for the database it owns; \texttt{CLINICAL\_DB\_*}, \texttt{ML\_DB\_*} on \texttt{analytics} for read-only credentials on upstream databases; \texttt{RABBITMQ\_URL} for the broker; \texttt{JWT\_SECRET\_KEY}, \texttt{JWT\_ALGORITHM}, \texttt{JWT\_EXPIRE\_MINUTES} on \texttt{auth} and the services that verify tokens; \texttt{GEMINI\_API\_KEY}, \texttt{GPT\_KEY} on \texttt{analytics} and \texttt{ml-batch} for the language model providers.

A \texttt{.env} file at the repository root provides defaults for development; in production the same variables are injected by the orchestrator (Compose, Kubernetes secret, or systemd drop-in). No credential is baked into a Docker image, so the same images can be promoted across environments without rebuilds.
